% ============================================================
%  Introduction Générale
% ============================================================
\chapter*{Introduction Générale}
\addcontentsline{toc}{chapter}{Introduction Générale}
\markboth{Introduction Générale}{}

Les organisations produisent et accumulent aujourd'hui des volumes de documents
internes sans précédent : rapports d'activité, procédures, manuels techniques,
contrats, factures, bons de commande, bases de connaissances. Ce patrimoine
informationnel constitue une ressource stratégique, mais il reste largement
sous-exploité. L'information existe, elle est archivée, mais elle demeure
difficilement accessible : retrouver un chiffre précis ou une procédure
particulière suppose de savoir dans quel document chercher, puis de parcourir
manuellement ce document. Le coût de cette recherche croît linéairement avec la
taille du corpus, jusqu'à devenir prohibitif.

Les moteurs de recherche documentaires classiques, fondés sur la correspondance
de mots-clés, n'apportent qu'une réponse partielle à ce problème. Ils retournent
une liste de documents, non une réponse : l'utilisateur doit encore lire et
synthétiser. Surtout, ils échouent dès que la formulation de la question ne
reprend pas le vocabulaire exact du document --- une requête sur les
« congés » ne trouvera pas un texte parlant d'« absences autorisées ».

L'émergence des grands modèles de langue (\eng{Large Language Models}, LLM) a
transformé cette situation. Ces modèles savent comprendre une question formulée
en langage naturel et produire une réponse rédigée. Employés seuls, ils
présentent toutefois deux limites disqualifiantes en contexte professionnel.
D'abord ils ne connaissent pas les documents internes de l'entreprise, qui ne
font pas partie de leurs données d'entraînement. Ensuite, et plus gravement, ils
\emph{hallucinent} : lorsqu'ils ignorent une réponse, ils en produisent une
plausible mais fausse, sans signaler leur incertitude. Un assistant interne qui
inventerait un montant de facture ou une clause contractuelle serait non
seulement inutile, mais dangereux.

L'architecture \emph{RAG} (\eng{Retrieval-Augmented Generation}) répond
précisément à ces deux limites. Son principe est d'ancrer la génération dans des
documents réels : plutôt que d'interroger le modèle sur ses connaissances
internes, on recherche d'abord les passages pertinents dans le corpus de
l'entreprise, puis on demande au modèle de répondre \emph{uniquement} à partir
de ces passages, en citant ses sources. La connaissance reste dans les
documents ; le modèle n'est plus qu'un moteur de reformulation contraint.

\section*{Objectifs du stage}

Ce stage, d'une durée de \internshipduration, a pour objet la conception, la
réalisation et le déploiement d'un système RAG complet permettant d'interroger
en langage naturel un corpus de documents d'entreprise. Cinq objectifs ont été
fixés :

\begin{enumerate}[label=\textbf{\arabic*.}]
  \item collecter, nettoyer et structurer les documents internes et mettre en
        place un pipeline d'indexation automatique ;
  \item générer des \eng{embeddings} sémantiques et les stocker dans une base de
        données vectorielle ;
  \item connecter la recherche sémantique à un LLM via un \eng{framework}
        d'orchestration afin de produire des réponses sourcées ;
  \item évaluer la fiabilité du système et réduire les hallucinations à l'aide
        d'un \eng{framework} d'évaluation dédié ;
  \item déployer l'assistant via une API et une interface web dans un
        environnement conteneurisé.
\end{enumerate}

Une contrainte transverse a structuré l'ensemble des choix techniques :
l'intégralité de la chaîne devait fonctionner \textbf{localement et sans coût
d'API}. Les documents internes d'une entreprise sont par nature confidentiels ;
les transmettre à un service tiers pour vectorisation ou génération soulève des
questions de conformité difficilement acceptables. Cette contrainte, loin d'être
une limitation, s'est révélée être un argument central en faveur de
l'architecture retenue.

\section*{Méthodologie et plan du rapport}

Le projet a été conduit selon la méthodologie \textbf{CRISP-DM}
(\eng{Cross-Industry Standard Process for Data Mining}), qui structure un projet
de science des données en six phases itératives. Le plan de ce rapport suit
fidèlement cette progression.

Le \textbf{chapitre 1} présente l'organisme d'accueil et détaille la
méthodologie CRISP-DM. Le \textbf{chapitre 2} formalise la problématique
métier, analyse les solutions existantes et justifie la solution retenue. Le
\textbf{chapitre 3} décrit le corpus documentaire, son exploration et sa
préparation. Le \textbf{chapitre 4} expose l'architecture technique du système
RAG : vectorisation, indexation, recherche sémantique et génération contrainte.
Le \textbf{chapitre 5} est consacré à l'évaluation de la fiabilité du système et
aux leviers mis en place contre les hallucinations. Le \textbf{chapitre 6}
présente le déploiement : API REST, interface web et conteneurisation. Une
conclusion générale dresse le bilan du travail réalisé et ouvre sur des
perspectives d'amélioration.
